\documentclass[12pt,a4paper]{article}
% Shared preamble for NLP course exercises — input this from each exercise file.

\usepackage{fontspec}
\setmainfont{Times New Roman}

\usepackage[a4paper, top=2.5cm, bottom=2.5cm, left=3cm, right=2.5cm]{geometry}

\usepackage{amsmath}
\usepackage{amssymb}
\usepackage{mathtools}
\usepackage{bm}

\usepackage{enumitem}
\usepackage{booktabs}
\usepackage{xcolor}
\usepackage{hyperref}
\usepackage{fancyhdr}

% Page style
\setlength{\headheight}{14pt}
\pagestyle{fancy}
\fancyhf{}
\lhead{\small\textit{Natural Language Processing}}
\rhead{\small Page \thepage}
\renewcommand{\headrulewidth}{0.4pt}

% Spacing
\setlength{\parindent}{0pt}
\setlength{\parskip}{6pt}

% Math operators
\DeclareMathOperator{\ReLU}{ReLU}


% ── Solution-specific helpers ─────────────────────────────────────────────────
\usepackage{mdframed}

% Light-gray box for final answers
\newmdenv[
  backgroundcolor = gray!12,
  linecolor       = gray!50,
  linewidth       = 0.8pt,
  innertopmargin  = 6pt,
  innerbottommargin = 6pt,
  innerleftmargin = 8pt,
  innerrightmargin = 8pt,
  skipabove       = 6pt,
  skipbelow       = 6pt,
]{answerbox}

% Step label: bold sans-serif marker
\newcommand{\step}[1]{\medskip\noindent\textbf{\textsf{Step #1.}}\enspace}

\begin{document}

% ── Title block ───────────────────────────────────────────────────────────────
\begin{center}
  {\Large\bfseries Natural Language Processing}\\[0.4em]
  {\large\bfseries Assignment 1 --- SOLUTION}\\[0.2em]
  {\large MLP for Binary Classification}\\[0.2em]
  {\normalsize Total: 10 points}
\end{center}

\noindent\rule{\linewidth}{0.6pt}
\vspace{0.5em}

% ═════════════════════════════════════════════════════════════════════════════
\section*{Question A \hfill{\normalfont\small(2 points)}}

\noindent\textit{Explain the role of the hidden layer, the ReLU activation,
and the Sigmoid output.}

\medskip

\textbf{Hidden layer.}
The hidden layer applies an affine (linear) transformation
\[
  \mathbf{a}^{(1)} = W^{(1)}\mathbf{x} + \mathbf{b}^{(1)},
\]
projecting the 2-dimensional input $\mathbf{x}$ into a new representation
space.
This transformation re-weights and mixes the input features, giving the
network the ability to detect patterns that are not directly visible in the
raw inputs.
Without a hidden layer the model degenerates to logistic regression, which can
only learn a \emph{linear} decision boundary.

\medskip

\textbf{ReLU activation.}
The Rectified Linear Unit (ReLU) is defined element-wise as
\[
  \ReLU(a) = \max(0,\, a).
\]
Its purpose is to introduce \emph{non-linearity}: without a non-linear
activation, a stack of linear layers collapses to a single linear
transformation, giving no benefit from depth.
ReLU specifically zeroes out any negative pre-activation, producing sparse
hidden representations.
Compared with sigmoid or tanh, ReLU does not saturate for large positive
inputs, which reduces the \emph{vanishing-gradient} problem during
backpropagation.

\medskip

\textbf{Sigmoid output.}
The sigmoid function
\[
  \sigma(z) = \frac{1}{1+e^{-z}}
\]
maps the unbounded scalar $z \in (-\infty,+\infty)$ to the open interval
$(0,1)$, making the output directly interpretable as the probability
$\hat{y} = P(y = 1 \mid \mathbf{x})$.
The standard decision rule is: predict class $1$ if $\hat{y} \ge 0.5$
(equivalently, $z \ge 0$), and class $0$ otherwise.

% ═════════════════════════════════════════════════════════════════════════════
\section*{Question B \hfill{\normalfont\small(6 points)}}

\noindent\textit{Perform the complete forward propagation.}

\bigskip
% ─────────────────────────────────────────────────────────────────────────────
\subsection*{(a) Compute $\mathbf{a}^{(1)}$ and $\mathbf{h}$
             \hfill{\normalfont\small(2 points)}}

\step{1} Apply the hidden-layer affine transformation.

\[
  \mathbf{a}^{(1)}
  = W^{(1)}\mathbf{x} + \mathbf{b}^{(1)}
  = \begin{bmatrix} 1 & -1 \\ 0.5 & 0.5 \end{bmatrix}
    \begin{bmatrix} 1 \\ 2 \end{bmatrix}
  + \begin{bmatrix} 0 \\ 0 \end{bmatrix}.
\]

Carry out the matrix--vector product row by row:

\[
  \begin{aligned}
    a^{(1)}_1 &= (1)(1) + (-1)(2) + 0 = 1 - 2 = -1,\\[4pt]
    a^{(1)}_2 &= (0.5)(1) + (0.5)(2) + 0 = 0.5 + 1.0 = 1.5.
  \end{aligned}
\]

\[
  \therefore\quad
  \mathbf{a}^{(1)} = \begin{bmatrix} -1 \\ 1.5 \end{bmatrix}.
\]

\step{2} Apply the ReLU activation element-wise
         $\bigl(\ReLU(a) = \max(0,a)\bigr)$.

\[
  \begin{aligned}
    h_1 &= \max(0,\,-1) = 0,       & &\text{(negative value is suppressed)}\\
    h_2 &= \max(0,\;1.5) = 1.5.   & &\text{(positive value passes through)}
  \end{aligned}
\]

\begin{answerbox}
\[
  \mathbf{a}^{(1)} = \begin{bmatrix} -1 \\ 1.5 \end{bmatrix}, \qquad
  \mathbf{h} = \ReLU\!\left(\mathbf{a}^{(1)}\right) = \begin{bmatrix} 0 \\ 1.5 \end{bmatrix}.
\]
\end{answerbox}

\noindent\textit{Observation:} Neuron 1 of the hidden layer is
\emph{inactive} (output $= 0$) because its pre-activation is negative.
Only neuron 2 carries information forward.

\bigskip
% ─────────────────────────────────────────────────────────────────────────────
\subsection*{(b) Compute $z$ and $\hat{y}$
             \hfill{\normalfont\small(3 points)}}

\step{1} Apply the output-layer affine transformation.

\[
  z = W^{(2)}\mathbf{h} + b^{(2)}
    = \begin{bmatrix} 1 & 1 \end{bmatrix}
      \begin{bmatrix} 0 \\ 1.5 \end{bmatrix}
    + 0
    = (1)(0) + (1)(1.5)
    = 1.5.
\]

\step{2} Compute $e^{-1.5}$ numerically.

Using $e^{-1.5} = e^{-1} \times e^{-0.5}$:
\[
  e^{-1} \approx 0.3679, \qquad
  e^{-0.5} \approx 0.6065, \qquad
  e^{-1.5} \approx 0.3679 \times 0.6065 \approx 0.2231.
\]

\step{3} Apply the sigmoid function.

\[
  \hat{y} = \sigma(1.5)
           = \frac{1}{1 + e^{-1.5}}
           = \frac{1}{1 + 0.2231}
           = \frac{1}{1.2231}
           \approx 0.8176.
\]

\begin{answerbox}
\[
  z = 1.5, \qquad \hat{y} = \sigma(1.5) \approx \boxed{0.8176}.
\]
\end{answerbox}

\noindent\textit{Interpretation:} The model estimates an $81.76\%$ probability
that the sample belongs to class $1$.

\bigskip
% ─────────────────────────────────────────────────────────────────────────────
\subsection*{(c) Predicted class \hfill{\normalfont\small(1 point)}}

Comparing $\hat{y}$ with the decision threshold $0.5$:
\[
  \hat{y} \approx 0.8176 > 0.5.
\]

\begin{answerbox}
The model predicts class $\hat{c} = \mathbf{1}$.
\end{answerbox}

% ═════════════════════════════════════════════════════════════════════════════
\section*{Question C \hfill{\normalfont\small(2 points)}}

\noindent\textit{Given $y = 1$, compute the Binary Cross-Entropy loss.}

\medskip

\step{1} Substitute $y = 1$ into the BCE formula.

\[
  \mathcal{L}(1,\,\hat{y})
  = -\bigl[\underbrace{(1)\log\hat{y}}_{y\,\log\hat{y}}
          + \underbrace{(1-1)\log(1-\hat{y})}_{=\;0}\bigr]
  = -\log\hat{y}.
\]

When the true label is 1, the loss reduces to the negative log of the
predicted probability.  The $(1-y)\log(1-\hat{y})$ term vanishes because
$1 - y = 0$.

\step{2} Substitute $\hat{y} = \sigma(z)$ and simplify using the
         log-sigmoid identity.

A useful identity avoids rounding errors from computing $\log \hat{y}$
indirectly:
\[
  -\log\sigma(z) = \log\!\left(1 + e^{-z}\right).
\]

Applying this with $z = 1.5$:
\[
  \mathcal{L}
  = \log\!\left(1 + e^{-1.5}\right)
  = \log(1 + 0.2231)
  = \log(1.2231).
\]

\step{3} Evaluate $\ln(1.2231)$ numerically.

\[
  \ln(1.2231)
  = \ln\!\left(\frac{1}{0.8176}\right)
  = -\ln(0.8176).
\]

Using the series expansion $\ln(1+x) \approx x - \tfrac{x^2}{2} +
\tfrac{x^3}{3} - \cdots$ with $x = 0.2231$:
\[
  \ln(1.2231) \approx 0.2231 - \frac{0.2231^2}{2} + \frac{0.2231^3}{3}
             \approx 0.2231 - 0.0249 + 0.0037
             \approx 0.2014 - 0.0001
             \approx 0.2014.
\]

\begin{answerbox}
\[
  \mathcal{L}(1,\,\hat{y}) = \ln(1 + e^{-1.5}) \approx \boxed{0.2014}.
\]
\end{answerbox}

\noindent\textit{Interpretation:} The loss is close to zero, reflecting that
the model assigns a high probability (${\approx}81.76\%$) to the correct class.
A perfect prediction ($\hat{y} = 1$) would give $\mathcal{L} = 0$; a random
guess ($\hat{y} = 0.5$) would give $\mathcal{L} = \ln 2 \approx 0.693$.

\end{document}
